%! AP Statistics — Unit 2, Lesson 2.2 — Group Activity KEY
\documentclass[10pt]{article}
\usepackage{apstats-article}
\usepackage{apstats-key}

%=============================================================================
\begin{document}

\pageheader{Unit 2, Lesson 2.2}{Group Activity --- Answer Key}
\begin{tabularx}{\linewidth}{X X X}
Name:\;\ans{Key} & Partner:\; & Period:\;
\end{tabularx}

\vspace{0.1in}

\begin{tcolorbox}[colback=sky, colframe=navy, boxrule=0.8pt, arc=2mm,
    left=3mm, right=3mm, top=2mm, bottom=2mm]
\small
Work with your group on the tier assigned by your teacher.
For each table: (1) compute conditional relative frequencies, (2) decide whether
the two variables are \textbf{associated}, and (3) justify using specific values.
\end{tcolorbox}

\vspace{0.12in}

% ════════════════════════════════════════════════════════════════════════════
% TIER R KEY
% ════════════════════════════════════════════════════════════════════════════
\begin{tcolorbox}[colback=navylight!30, colframe=navy, boxrule=0.9pt, arc=2mm,
    left=3mm, right=3mm, top=2mm, bottom=2mm,
    title={\bfseries Tier R --- Remediate}, fonttitle=\bfseries\small, halign title=center]
\small

\textbf{Context:} 120 students --- after-school job $\times$ school sport.

\vspace{8pt}
\renewcommand{\arraystretch}{1.8}
\begin{tabularx}{\linewidth}{@{} l C{2.4cm} C{2.4cm} C{2.2cm} @{}}
\rowcolor{sky}
             & \textbf{Sport: Yes} & \textbf{Sport: No} & \textbf{Row Total} \\
\hline
Job: Yes & 12 & 28 & 40 \\
Job: No  & 48 & 32 & 80 \\
\hline
Column Total & 60 & 60 & 120 \\
\end{tabularx}

\vspace{10pt}
\textbf{1.} Conditional relative frequencies by job group:

\vspace{6pt}
\renewcommand{\arraystretch}{1.8}
\begin{tabularx}{\linewidth}{@{} l C{2.6cm} C{2.6cm} C{1.8cm} @{}}
\rowcolor{sky}
             & \textbf{Sport: Yes} & \textbf{Sport: No} & \textbf{Total} \\
\hline
Job: Yes & \ans{$12/40=30\%$} & \ans{$28/40=70\%$} & 100\% \\
Job: No  & \ans{$48/80=60\%$} & \ans{$32/80=40\%$} & 100\% \\
\end{tabularx}

\vspace{10pt}
\textbf{2.} \ans{\textbf{Yes}, there is evidence of an association. Among students with
a job, only 30\% play a sport; among students without a job, 60\% play a sport.
The conditional distributions differ substantially, suggesting that having a job is
associated with lower participation in school sports.}

\end{tcolorbox}

\vspace{0.14in}

% ════════════════════════════════════════════════════════════════════════════
% TIER A KEY
% ════════════════════════════════════════════════════════════════════════════
\begin{tcolorbox}[colback=greenbg, colframe=greenacc, boxrule=0.9pt, arc=2mm,
    left=3mm, right=3mm, top=2mm, bottom=2mm,
    title={\bfseries Tier A --- Approaching Proficiency}, fonttitle=\bfseries\small, halign title=center]
\small

\textbf{Context:} 180 students --- class year $\times$ movie genre.

\vspace{6pt}
\renewcommand{\arraystretch}{1.8}
\begin{tabularx}{\linewidth}{@{} l C{2.2cm} C{2.2cm} C{2.2cm} C{2.2cm} @{}}
\rowcolor{sky}
                    & \textbf{Action} & \textbf{Comedy} & \textbf{Drama} & \textbf{Row Total} \\
\hline
Freshman   & 36 & 54 & 30 & \ans{120} \\
Sophomore  & 30 & 18 & 12 & \ans{60} \\
\hline
Col.\ Total& \ans{66} & \ans{72} & \ans{42} & \ans{180} \\
\end{tabularx}

\vspace{10pt}
\textbf{1.} Conditional relative frequencies:

\vspace{6pt}
\renewcommand{\arraystretch}{1.8}
\begin{tabularx}{\linewidth}{@{} l C{2.4cm} C{2.4cm} C{2.4cm} C{1.6cm} @{}}
\rowcolor{sky}
                    & \textbf{Action} & \textbf{Comedy} & \textbf{Drama} & \textbf{Total} \\
\hline
Freshman  & \ans{30\%} & \ans{45\%} & \ans{25\%} & 100\% \\
Sophomore & \ans{50\%} & \ans{30\%} & \ans{20\%} & 100\% \\
\end{tabularx}

\vspace{10pt}
\textbf{2.} Segmented bar chart: two bars (Freshman / Sophomore), each at 100\%, divided
into Action (bottom), Comedy (middle), Drama (top). Freshman bars: 30\% / 45\% / 25\%;
Sophomore bars: 50\% / 30\% / 20\%.

\vspace{10pt}
\textbf{3.} \ans{There is evidence of an association between class year and genre preference.
Among freshmen, 45\% prefer comedy, compared to only 30\% among sophomores; sophomores
prefer action at a higher rate (50\% vs.\ 30\%). The conditional distributions differ
noticeably, so the two variables appear to be associated.}

\end{tcolorbox}

\vspace{0.14in}

% ════════════════════════════════════════════════════════════════════════════
% TIER E KEY
% ════════════════════════════════════════════════════════════════════════════
\begin{tcolorbox}[colback=redbg, colframe=redacc, boxrule=0.9pt, arc=2mm,
    left=3mm, right=3mm, top=2mm, bottom=2mm,
    title={\bfseries Tier E --- Extension}, fonttitle=\bfseries\small, halign title=center]
\small

\textbf{Context:} 300 adults --- commute method $\times$ stress level (joint relative frequencies given).

\vspace{6pt}
\renewcommand{\arraystretch}{1.8}
\begin{tabularx}{\linewidth}{@{} l C{2.4cm} C{2.4cm} C{2.4cm} @{}}
\rowcolor{sky}
                  & \textbf{Stress: Low} & \textbf{Stress: High} & \textbf{Row Total} \\
\hline
Car     & 0.20 & 0.30 & \ans{0.50} \\
Transit & 0.15 & 0.15 & \ans{0.30} \\
Walk    & 0.14 & 0.06 & \ans{0.20} \\
\hline
Col.\ Total & \ans{0.49} & \ans{0.51} & \ans{1.00} \\
\end{tabularx}

\vspace{10pt}
\textbf{1.} \ans{$0.20+0.30+0.15+0.15+0.14+0.06 = 1.00$ \checkmark}

\vspace{8pt}
\textbf{2.} Conditional P(High stress $|$ commute):

\vspace{6pt}
\renewcommand{\arraystretch}{1.8}
\begin{tabularx}{\linewidth}{@{} l C{3.4cm} @{}}
\rowcolor{sky} \textbf{Commute} & \textbf{P(High stress $|$ commute)} \\
\hline
Car     & \ans{$0.30/0.50 = 60\%$} \\
Transit & \ans{$0.15/0.30 = 50\%$} \\
Walk    & \ans{$0.06/0.20 = 30\%$} \\
\end{tabularx}

\vspace{8pt}
\textbf{3.} \ans{There is evidence of an association between commute method and stress level.
Car commuters have the highest conditional stress rate (60\%), while walkers have the
lowest (30\%). The large difference across groups indicates an association.}

\vspace{6pt}
\textbf{4.} \ans{Conditioning on stress level instead: P(Car $|$ Low) $= 0.20/0.49 \approx 41\%$;
P(Car $|$ High) $= 0.30/0.51 \approx 59\%$. The association conclusion does
\textbf{not} change --- high-stress adults use cars at a higher rate than low-stress adults.
Association is symmetric: if $X$ is associated with $Y$, then $Y$ is associated with $X$.
The direction of conditioning changes the framing but not whether an association exists.}

\end{tcolorbox}

\begin{teachernote}
\textbf{Tier R}: Watch for students dividing by 120 (grand total) instead of 40 or 80 (row totals).
\textbf{Tier A}: Row totals are unequal (120 and 60) --- a key teaching moment for why raw counts are misleading and conditional relative frequencies are needed.
\textbf{Tier E}: Students must recognize that joint relative frequencies don't sum to 1.00 per row; they must still divide by row totals. Stretch item (Q4) reinforces that association is symmetric.
\end{teachernote}

\end{document}
